% !TEX program = xelatex
\IfFileExists{main.tex}{\documentclass[main.tex]{subfiles}}{
  \IfFileExists{../../main.tex}{\documentclass[../../main.tex]{subfiles}}{
    \IfFileExists{../main.tex}{\documentclass[../main.tex]{subfiles}}{
      \documentclass[Book\space NL/main.tex]{subfiles}
    }
  }
}
\begin{document}
\ifSubfilesClassLoaded{
  \ifcsname appendixchapterparent\endcsname
  \else
    \standaloneappendixchaptersetup{C}{3}
  \fi
}{}
\chapter{Medical Mycology}
\labch{appendix-c3-medical-mycology}
\margintoc{}

Fungi are defined as heterotrophic organisms of variable morphology (they can appear in yeast and mold forms) which perform extracellular digestion. They typically possess chitin, glucan, and mannan cell-wall components, with an ergosterol-containing cell membrane. Fungi reproduce through spores or budding. Fungi can be further defined by their mechanism of feeding: saprophytic fungi feed from dead organisms, while parasitic fungi absorb nutrients from living tissues.

Yeasts are described as unicellular fungi\footnote{Though they may present pseudohyphae, which are individual cells linked together in hyphae-like morphology.} while moulds are multicellular fungi that grow into cord-like structures. Typically, molds are aerobic organisms, while yeasts are facultative anaerobes. Some fungi are able to exhibit a mixture of these characteristics depending on their environmental conditions. These fungi are considered dimorphic.

Fungal hyphae may be distinguished by the presence or absence of septa (septate or aseptate hyphae) and by colour (moniliaceous = colourless; dematiaceous = pigmented).

\marginnote{
  Key Terms:
  \begin{itemize}
    \item Hypha (hyphae) ---  tubular cell walls surrounding the plasma membrane and cytoplasm of cells.
    \item Mycelium (mycelia) --- interwoven massive form of hyphae.
  \end{itemize}
}

\textbf{Fungal states}
\begin{itemize}
  \item Anamorph: the asexual reproductive state.
  \item Teleomorph: the sexual reproductive state.
  \item Holomorph: the whole fungus, including both anamorph and teleomorph states.
\end{itemize}

Reproductive processes of fungi:
\begin{itemize}
  \item Sexual reproduction --- plasmogamy (fusion of cytoplasm), followed by karyogamy (fusion of nuclei), meiosis, spore formation, and germination.
  \item Asexual reproduction --- mitosis followed by germination; conidia produce mitospores.
\end{itemize}

Spores are similar to seeds with high resilience and the ability to germinate into fungi. They may be produced by either sexual or asexual means (see above). Examples: ascospores, basidiospores and zygospores.

Conidia: Asexual reproductive structures (mitospores) produced either from the transformation of a vegetative yeast or hyphal cell or from a specialized conidiogenous cell. Conidia may be formed on specialized hyphae, termed conidiophores.
\begin{itemize}
  \item Based on size, microconidia are small, whereas macroconidia are large or multicellular.
  \item Arthroconidia (arthrospores): conidia that result from the fragmentation of hyphal cells.
  \item Blastoconidia (blastospores): conidia formed through a budding process (e.g., in yeasts).
  \item Chlamydoconidia (chlamydospores): large, thick-walled, usually spherical conidia produced from terminal or intercalary hyphal cells.
  \item Phialoconidia: conidia produced by a vase-shaped conidiogenous cell called a phialide (e.g., \textit{Aspergillus fumigatus}).
  \item Sporangiospores: asexual structures characteristic of the order Mucorales; they are mitotic spores produced within an enclosed sporangium, often supported by a sporangiophore.
\end{itemize}

\section{Virulence Factors of Fungi}
1. Thermotolerance
2. Dimorphism --- help withstand host defenses
3. Cell wall components and capsule --- help withstand host defenses. Includes alpha 1--3-glucan
4. Adhesion molecules
5. Enzyme production
• Catalase
• Superoxide dismutase
• Extracellular proteinases
  - Adherence
  - Survival of the pathogen on mucosal surfaces
  - Invasion of host tissues
  - Digestion of immunoglobulins
• Keratinase
• Elastase
• Collagenase
6. Miscellaneous
Antigenic variation
– Phenotypic switching
• Thigmotropism (bend on contact)
• Toxin
• Intracellular survival
• Drug resistance

\chapter{Fungal Identification}
\margintoc{}

Specimens collected from the site of infection. Blood is use for systemic infect. Use of sterile technique and promptly transport sample for testing. Includes conventional, histopathologic, molecular and biochemical method.

\section{Laboratory Diagnosis of Fungal Infections}

\textbf{Direct microscopic examination}\index{microscopy of fungi} is used to visualize hyphae, yeast cells, or spores from clinical specimens, allowing for the immediate start of treatment. Key stains include potassium hydroxide (of skin/tissue scrapings to lyse tissue), lactophenol cotton blue (lactic acid for preserving and cotton blue for staining) and Nigrosin (capsule negative stain --- \emph{Cryptococcus}, esp, \emph{C. neoformans}). Tissue stains may also be performed, including Gram staining as well. Giemsa stain can be used for \emph{Histoplasma capsulatum}.

Selected stains:
\begin{itemize}
  \item \textbf{Calcofluor white stain}\index{Calcofluor white} binds to chitin and cellulose components of the fungal cell wall, allowing fungus to be fluorescently labelled.
  \item \textbf{Periodic acid Schiff}\index{periodic acid--Schiff stain} stains the fungal cell wall in tissue samples. This is by the periodic acid oxidizes polysaccharides in fungal cell walls to aldehydes, which then react with the Schiff reagent to produce a magenta color.
  \item \textbf{Gomori methenamine silver}\index{Gomori methenamine silver} stains fungi in biopsies by binding the fungal cell wall to appear black coloured.
  \item \textbf{Mucicarmine stain}\index{Mucicarmine stain} stains \emph{mucin-producing fungi}, esp. Cryptococcus by binding acidic mucopolysaccharides in cell capsule to rose pink.
\end{itemize}

Cultures may also be used in fungal identification. Fungi can be rapidly growing or slow growing, as well as fastidious. Culture media include:
\textbf{Enriched media:}
\begin{itemize}
  \item Sabouraud dextrose (SDA) agar
  \item Brain-heart infusion agar
\end{itemize}

\textbf{Selective media:}
\begin{itemize}
  \item Chloramphenicol and cycloheximide --- suppress bacterial and contaminating fungal growth
  \item CHROMagar --- selective identification of \emph{C. albicans, C. tropicalis, C. krusei}.
  \item Sunflower/bird seed --- for \emph{Cryptococcus}.
\end{itemize}

\section{Fungal Identification}

Observation of colony characteristics provides identification clues. These include colony morphology and texture, colour, growth rate, and topography.

Additional tests include the germ-tube test and chlamydospore formation for \emph{C. albicans}. Biochemical tests are used for many other species. Immunological methods include latex agglutination tests.



\section{High-yield fungal pathogen reference}
\label{sec:high-yield-fungal-pathogens}

The following table consolidates the fungal-organism notes in the accompanying review workbook.

\pagelayout{wide}
\begingroup
\footnotesize
\setlength{\tabcolsep}{3pt}
\renewcommand{\arraystretch}{1.15}
\begin{longtable}{@{}>{\RaggedRight\arraybackslash}p{0.19\linewidth}>{\RaggedRight\arraybackslash}p{0.30\linewidth}>{\RaggedRight\arraybackslash}p{0.24\linewidth}>{\RaggedRight\arraybackslash}p{0.21\linewidth}@{}}
\caption[Selected fungal pathogens]{Selected fungal and fungus-like pathogens: clinical features, diagnostic clues, and management notes.}\label{tab:fungal-pathogen-summary}\\
\toprule
\textbf{Organism / category} & \textbf{Clinical features / key characteristics} & \textbf{Diagnosis} & \textbf{Treatment / notes} \\
\midrule
\endfirsthead
\caption[]{Selected fungal pathogens (continued).}\\
\toprule
\textbf{Organism / category} & \textbf{Clinical features / key characteristics} & \textbf{Diagnosis} & \textbf{Treatment / notes} \\
\midrule
\endhead
\midrule
\multicolumn{4}{r}{\footnotesize Continued on next page}\\
\endfoot
\bottomrule
\endlastfoot
\textbf{Superficial mycoses} & & & \\
\textit{Malassezia} spp. (pityriasis versicolor / tinea versicolor) & Infection of the stratum corneum, usually on oily skin; hypo- or hyperpigmented patches with little inflammation. The organism is lipophilic and lipolytic. & KOH microscopy may show short hyphae and yeast cells; Wood lamp may fluoresce. & Topical azoles (e.g., ketoconazole), selenium sulfide, or zinc pyrithione.\\
\textit{Hortaea werneckii} (tinea nigra palmaris) & Halotolerant fungus causing a brown-to-black superficial macule, usually on the palm or sole. & KOH microscopy and fungal culture. & Topical antifungal treatment; salicylic acid may be used as a keratolytic adjunct.\\
Piedra: \textit{Piedraia hortae} / \textit{Trichosporon} spp. & Hair-shaft infection; black piedra is caused by \textit{Piedraia hortae}, while white piedra is associated with \textit{Trichosporon} spp. & Examine hair shafts and nodules by microscopy and culture. & Hair trimming or shaving and topical antifungal shampoo may be used.\\
\textit{Trichosporon} spp. & White piedra, onychomycosis, otomycosis, and, in susceptible hosts, invasive or disseminated trichosporonosis. \textit{T. asahii} may form barrel-shaped arthroconidia. & Culture with species identification; assess for cross-reaction with cryptococcal antigen tests. & Management depends on the infection site and susceptibility; invasive disease requires specialist treatment.\\
\textbf{Cutaneous / dermatophyte infections} & & & \\
Dermatophytes: \textit{Trichophyton}, \textit{Microsporum}, \textit{Epidermophyton} & Invade keratinised skin, hair, and nails. \textit{Trichophyton}: pencil-shaped macroconidia; \textit{Microsporum}: spindle-shaped macroconidia; \textit{Epidermophyton}: club-shaped macroconidia. & KOH microscopy and culture on Sabouraud dextrose agar or dermatophyte test medium. Wood lamp may help with selected species. Hair-perforation and urease tests may assist speciation. & Topical antifungals for limited skin disease; oral therapy is usually required for tinea capitis, extensive disease, or nail involvement.\\
Tinea pedis / tinea manuum & Dermatophyte infection of the feet / hands. Entry may occur through skin abrasions. & KOH microscopy and culture when the diagnosis is uncertain or disease is recurrent. & Keep affected areas clean and dry; use topical antifungal therapy.\\
Tinea corporis / tinea cruris & Ring-shaped lesions on the trunk or limbs / groin and adjacent skin ("jock itch"). & KOH microscopy of skin scrapings; culture if atypical or refractory. & Topical antifungal treatment; assess for extensive disease or steroid-modified tinea.\\
Tinea capitis & Infection of scalp hair; kerion and favus may cause marked inflammation. & KOH microscopy, fungal culture, and species identification. & Requires systemic antifungal therapy; assess for inflammatory kerion and scarring.\\
Tinea barbae / tinea faciei & Dermatophyte infection of the beard area / face. & KOH microscopy and fungal culture. & Treatment depends on extent and hair-follicle involvement; systemic therapy may be required.\\
Tinea unguium (onychomycosis) & Dermatophyte infection of fingernails or toenails; non-dermatophyte moulds and yeasts can also cause onychomycosis. & Nail clipping or scraping for microscopy, histopathology, and culture. & Confirm the aetiology before prolonged systemic antifungal treatment.\\
Tinea imbricata / tinea incognito & Tinea imbricata is a superficial mycosis caused by \textit{Trichophyton concentricum}. Tinea incognito has atypical features after corticosteroid or immunosuppressant use. & KOH microscopy and culture; review prior topical steroid use. & Treat the underlying dermatophyte infection and avoid inappropriate topical corticosteroid monotherapy.\\
Deep dermatophytosis & Deep dermatophytosis, disseminated dermatophytosis, or dermatophyte abscess may occur, particularly in people with chronic illness or impaired immunity. & Histopathology, fungal culture, and species identification. & Specialist management; systemic antifungal treatment is generally required.\\
\textbf{Yeasts and opportunistic fungi} & & & \\
\textit{Candida albicans} & Causes mucocutaneous and invasive candidiasis. Risk factors include immunosuppression, broad-spectrum antibiotics, parenteral nutrition, catheters, and other healthcare exposures. Virulence factors include adhesins, biofilm formation, dimorphism, thigmotropism, proteases, phospholipases, and phenotypic switching. & Direct examination may show budding yeast and pseudohyphae; Gram stain, wet mount, and KOH preparation. Culture, germ-tube test, chlamydospore formation, CHROMagar, MALDI-TOF, and molecular tests may be used. & Treatment depends on site, severity, species, and susceptibility; systemic disease requires prompt treatment and assessment of intravascular devices.\\
Non-\textit{albicans Candida} spp. & Includes \textit{C. tropicalis}, \textit{C. glabrata}, \textit{C. krusei}, and \textit{C. parapsilosis}; morphology and antifungal susceptibility vary by species. & Culture with species identification and antifungal susceptibility testing. & Select therapy according to species and susceptibility; some species have reduced or intrinsic susceptibility to commonly used azoles.\\
\textit{Candida auris} & Healthcare-associated yeast that can cause outbreaks and may be multidrug-resistant. & Accurate identification by validated methods (e.g., MALDI-TOF or sequencing) and susceptibility testing. & Infection-control measures are essential; treatment should be susceptibility-guided and managed with specialist input.\\
\textit{Fusarium}, \textit{Aspergillus}, and \textit{Candida} spp. (fungal keratitis) & Important causes of corneal infection. & Corneal scraping for direct microscopy and culture; urgent ophthalmic assessment. & Specialist ophthalmic treatment; topical antifungal therapy is selected according to the organism.\\
\textit{Cryptococcus neoformans} and \textit{C. gattii} & Encapsulated yeasts causing pulmonary, CNS, cutaneous, or disseminated disease. \textit{C. gattii} can also infect immunocompetent people. & Cryptococcal antigen testing; India ink, mucicarmine, GMS, culture, urease, and phenol-oxidase testing may support diagnosis. & CNS disease generally requires induction therapy with amphotericin B plus flucytosine, followed by an oral azole; specialist management is required.\\
\textit{Pneumocystis jirovecii} & Causes pneumonia in immunocompromised people; may present with fever, dry cough, progressive dyspnoea, hypoxia, and ground-glass opacities. & Respiratory specimens, including BAL, for staining and PCR. & Trimethoprim--sulfamethoxazole is first-line treatment and is also used for prophylaxis in selected high-risk groups.\\
\textit{Rhodotorula} spp. & Pigmented yeasts associated with bloodstream infection, particularly in people with central venous catheters or immunosuppression. & Blood culture with species identification and susceptibility testing. & Assess for catheter-related infection; treatment is susceptibility-guided.\\
\textit{Saccharomyces cerevisiae} & Baker's and brewer's yeast; can occasionally cause opportunistic infection, particularly in vulnerable patients. & Culture and clinical correlation. & Manage as an opportunistic yeast infection; review probiotic exposure and intravascular devices where relevant.\\
\textit{Scytalidium} / \textit{Neoscytalidium} spp. and \textit{Geotrichum} spp. & May cause skin or nail infection; some species can mimic oral candidiasis. & Direct microscopy and fungal culture with species identification. & Treatment varies by organism and site; confirm the pathogen before selecting therapy.\\
\textbf{Other fungal or fungus-like pathogens} & & & \\
\textit{Pythium insidiosum} (oomycete) & Water-associated, fungus-like organism causing ocular, cutaneous, vascular, or other forms of pythiosis; it is not a true fungus. Zoospores are important in its life cycle. & Corneal scraping or tissue from the affected site; microscopy, culture, histopathology, serology, and molecular testing. & Management often requires specialist care and surgical intervention; response to antifungals is variable.\\
Microsporidia & Obligate intracellular organisms; historically studied with fungi, but not yeasts or moulds. & Tissue or stool microscopy with appropriate stains and molecular testing, selected according to the affected site. & Restore immune function where possible and use species-guided treatment.\\
\textit{Sporothrix schenckii} complex (sporotrichosis) & "Rose gardener's disease"; associated with soil and plant material. Papules may enlarge and ulcerate, with nodules spreading along lymphatics. & Culture; tissue may show cigar-shaped yeast or asteroid bodies. & Itraconazole is commonly used for lymphocutaneous disease; severe or disseminated disease needs specialist treatment.\\
Chromoblastomycosis agents & Chronic verrucous or cauliflower-like cutaneous lesions with epidermal hyperplasia. & Histopathology may show muriform (sclerotic) bodies; culture supports identification. & Prolonged systemic antifungal treatment, sometimes combined with surgery.\\
Mycetoma agents (eumycetoma / actinomycetoma) & Chronic subcutaneous infection with swelling, draining sinuses, and discharge containing grains; commonly affects the foot ("Madura foot"). & Examine grains and tissue; microscopy, culture, and histopathology distinguish fungal eumycetoma from bacterial actinomycetoma. & Treatment differs substantially between eumycetoma and actinomycetoma; establish the cause before therapy.\\
\end{longtable}
\endgroup
\pagelayout{margin}

\chapter{Antifungal agents}
\margintoc{}



\end{document}